% ============================================================
% estilo.tex — Plantilla visual del libro guía
% Teoría Matemática de la Luz (H. Poincaré) — edición en español
% ============================================================

% ---------- Geometría de página: más texto por página ----------
\usepackage[
  paperwidth=17cm,paperheight=24cm,
  inner=2cm,outer=1.8cm,top=2.2cm,bottom=2.4cm,
  headheight=35pt,includeheadfoot
]{geometry}

% ---------- Tipografía y codificación ----------
\usepackage[spanish,es-noindentfirst]{babel}
\usepackage[utf8]{inputenc}
\usepackage[T1]{fontenc}
\usepackage{newpxtext,newpxmath} % serif de corte clásico, afín al original de 1889
\usepackage{microtype}

% ---------- Matemáticas ----------
% newpxmath ya redefine \Bbbk y \openbox; se evita el choque con
% amssymb/amsthm (no se usan entornos de teorema: se usan las cajas
% tcolorbox definidas más abajo).
\usepackage{amsmath}
\let\Bbbk\relax
\usepackage{amssymb}
\usepackage{mathtools}

% ---------- Color: paleta del libro ----------
\usepackage{xcolor}
\definecolor{azulpoincare}{HTML}{14315D}      % azul profundo — títulos, estructura
\definecolor{azulclaro}{HTML}{E8EEF5}         % azul muy claro — fondo de cajas de definición
\definecolor{rojopoincare}{HTML}{9E1B32}      % rojo — acentos, cálculos, numeración
\definecolor{rojoclaro}{HTML}{F7E9EA}         % rojo muy claro — fondo de cajas de cálculo
\definecolor{grisnota}{HTML}{4A4A4A}          % gris — notas históricas

% ---------- Figuras y cajas ----------
\usepackage{graphicx}
\usepackage[most]{tcolorbox}
\tcbuselibrary{skins,breakable}

% ---------- Encabezados ----------
\usepackage{fancyhdr}
\usepackage{titlesec}
\usepackage{enumitem}
\usepackage{hyperref}
\hypersetup{
  colorlinks=true,
  linkcolor=azulpoincare,
  citecolor=rojopoincare,
  urlcolor=azulpoincare,
  pdftitle={Teoría Matemática de la Luz — Henri Poincaré},
}

% ============================================================
% Portadilla de capítulo — imita la composición tipográfica
% original ("CHAPITRE N" / título centrado, filete, blanco)
% pero en español y con la paleta azul/rojo
% ============================================================
\titleformat{\chapter}[display]
  {\normalfont\centering}
  {\color{rojopoincare}\Large\scshape Capítulo \thechapter}
  {12pt}
  {\color{azulpoincare}\Huge\bfseries\scshape}
\titlespacing*{\chapter}{0pt}{80pt}{40pt}

\renewcommand{\thesection}{\thechapter.\arabic{section}}

% ============================================================
% Cajas de estilo "libro guía"
% ============================================================

% Nota sobre los cálculos que extendemos: NO llevan caja, ni encabezado,
% ni pasos numerados, ni subtítulos en negrilla. Se escriben como
% continuación en prosa del propio texto, en el mismo registro que
% Poincaré, como una simple ampliación de lo que él deja implícito — no
% como un aparato pedagógico aparte. Ver CLAUDE.md.

% --- Enunciado importante / principio / resultado clave ---
% Para destacar afirmaciones centrales (principios, resultados
% culminantes), NO para sus demostraciones.
\newtcolorbox{enunciado}[1]{%
  breakable,
  enhanced,
  colback=azulclaro,
  colframe=azulpoincare,
  boxrule=0.6pt,
  left=8pt,right=8pt,top=6pt,bottom=6pt,
  fonttitle=\bfseries\color{white},
  coltitle=white,
  colbacktitle=azulpoincare,
  title={#1},
  attach boxed title to top left={xshift=8pt,yshift=-2mm},
  boxed title style={colback=azulpoincare,boxrule=0pt,sharp corners},
}

% --- Nota histórica / contexto ---
\newtcolorbox{notahistorica}[1][]{%
  breakable,
  enhanced,
  colback=white,
  colframe=grisnota,
  boxrule=0.5pt,
  left=8pt,right=8pt,top=6pt,bottom=6pt,
  fonttitle=\itshape\color{grisnota},
  coltitle=grisnota,
  title={Nota histórica#1},
  attach boxed title to top left={xshift=8pt,yshift=-2mm},
  boxed title style={colback=white,colframe=grisnota,boxrule=0.5pt,sharp corners},
}

% --- Definición / concepto clave ---
\newtcolorbox{definicion}[1][]{%
  breakable,
  enhanced,
  colback=azulclaro,
  colframe=azulpoincare,
  boxrule=0.6pt,
  left=8pt,right=8pt,top=6pt,bottom=6pt,
  fonttitle=\bfseries\color{white},
  coltitle=white,
  colbacktitle=azulpoincare,
  title={Definición#1},
  attach boxed title to top left={xshift=8pt,yshift=-2mm},
  boxed title style={colback=azulpoincare,boxrule=0pt,sharp corners},
}

% ============================================================
% Ecuaciones: numeración doble (libro guía + referencia original)
% Uso: \begin{ecnum}{III.4}{orig. (4)} ... \end{ecnum}
% ============================================================
\newenvironment{ecnum}[2]{%
  \gdef\ecnumA{#1}\gdef\ecnumB{#2}%
  \begin{equation}%
}{%
  \tag{\ecnumA\ \text{\footnotesize[\ecnumB]}}%
  \end{equation}%
}

% ---------- Encabezado de página (estilo "guía de lectura") ----------
\setlength{\headheight}{35pt}
\pagestyle{fancy}
\fancyhf{}
\fancyhead[LE,RO]{\thepage}
\fancyhead[RE]{\color{azulpoincare}\scshape\nouppercase{\leftmark}}
\fancyhead[LO]{\color{azulpoincare}\scshape\nouppercase{\rightmark}}
\renewcommand{\headrulewidth}{0.4pt}
\renewcommand{\headrule}{\hbox to\headwidth{\color{rojopoincare}\leaders\hrule height \headrulewidth\hfill}}
